\honest{Raising the Sanity Waterline}{Eliezer Yudkowsky}{2009}

Behind every dramatic failure, to paraphrase the Black Belt Bayesian, there is a larger and less dramatic failure that made it possible. If religion vanished tomorrow, the ``larger failures of sanity that made religion possible in the first place'' would remain. Religion is ``a direct problem'', worth fighting. It is also a dead canary in a coal mine: a symptom.

So I propose a thought experiment. Teach a course of general reasoning that never mentions religion, not even by hint, and five years later count how many of your students have lost their religion, compared with a control group. What would you teach so that ``religion goes underwater''? \nb{The post assumes that such a course would lower religious belief, and gives no evidence. A 2019 study by Dyer and Hall found that general research methods classes did not reduce students' beliefs in pseudoscience and the like, while critical thinking classes that addressed those beliefs directly did. Its abstract does not mention religion.} I list topics I have already written about that did not avoid religion but could: affective death spirals, which have plenty of examples outside the supernatural; cached thoughts, fake wisdom and the pressure to conform; evidence and Occam's Razor; why reason works causally; mysterious answers, with vitalism and phlogiston as historical examples; the absence of fundamental mental things, through the Mind Projection Fallacy applied to probability, then reductionism and cognitive science; the arts of actually updating on evidence, under some name other than ``Crisis of Faith''; Fun Theory, taught as a theory of utopian fiction; Joy in the Merely Real and naturalistic metaethics. Dark Side Epistemology ``would be hard'', but you might videotape the interrogation of a snake-oil salesman.

Now take a scientist who is religious, or who tosses around endorsements of ``spirituality''. ``We now know this person is not applying any technical, explicit understanding'' of evidence, Occam's Razor, curiosity-stoppers, updating on evidence, ``epistemology 101'' and ``self-honesty 201''. The list also includes how evidence and Occam's Razor follow from minds working as mapping engines, and so do not switch off when you talk about tooth fairies; thinking for yourself instead of repeating what you heard; and the general trends of science over three thousand years. \nb{In ``Outside the Laboratory'', which the post links here, Yudkowsky had offered this diagnosis as a personal view (``to me'') and with ``probably''. Here it is stated as knowledge, and again no real scientist's reasoning is examined.} An introduction to all of it would be one four-credit undergraduate course.

Yet Nobel laureates have not taken it: ``Richard Smalley if you're looking for a cheap shot, or Robert Aumann if you're looking for a scary shot.'' They cannot be isolated exceptions. Had their colleagues taken it, they would have corrected them, or pitied them too much to give them a Nobel. Could you, realistically, win a Nobel while advocating the existence of Santa Claus? \nb{Smalley won his prize in 1996. The creationist statements quoted in the comment the post links are from 2004 and 2005. For Aumann's own account of his religion, see ``Crisis of Faith''.} So the sanity waterline is ``really ridiculously low'', ``Even in the highest halls of science.'' Throw out the canary and the mine may stink less, but the waterline may not rise much.

Last, I back the neo-atheists. The harm done by religion is ``current and ongoing disaster''; I do not argue it. Fighting that harm comes before using religion as a canary. Even if Dawkins, Dennett, Harris and Hitchens won everywhere, ``the real work of rationalists will be only just beginning.''
